% Evaluation metrics (short section for paper)
% Usage: % Evaluation metrics (short section for paper)
% Usage: % Evaluation metrics (short section for paper)
% Usage: % Evaluation metrics (short section for paper)
% Usage: \input{paper_tables/evaluation_metrics.tex}

\subsection{Evaluation Metrics}
\label{sec:evaluation-metrics}

We evaluate all methods on a binary classification task:
Bug-Fix (minority class) versus Non-Bug-Fix.
Predictions are compared against human-annotated labels using
standard per-class and aggregate metrics derived from the
confusion matrix (true positives~TP, true negatives~TN,
false positives~FP, false negatives~FN).

For each class $c \in \{\text{Non-Bug-Fix}, \text{Bug-Fix}\}$, we
report \emph{precision} ($\mathrm{TP}_c / (\mathrm{TP}_c + \mathrm{FP}_c)$),
\emph{recall} ($\mathrm{TP}_c / (\mathrm{TP}_c + \mathrm{FN}_c)$), and
\emph{F1} (harmonic mean of precision and recall).
We additionally report \emph{accuracy}
($(\mathrm{TP} + \mathrm{TN}) / N$) and the
\emph{Matthews Correlation Coefficient} (MCC)~\cite{matthews1975comparison}:
\[
  \mathrm{MCC} =
  \frac{\mathrm{TP}\cdot\mathrm{TN} - \mathrm{FP}\cdot\mathrm{FN}}
       {\sqrt{(\mathrm{TP}+\mathrm{FP})(\mathrm{TP}+\mathrm{FN})
              (\mathrm{TN}+\mathrm{FP})(\mathrm{TN}+\mathrm{FN})}}.
\]

Because the dataset is imbalanced (599 Non-Bug-Fix vs.\ 199 Bug-Fix
commits), accuracy alone can be misleading: a classifier that predicts
Non-Bug-Fix for nearly every commit achieves high accuracy while failing
on the minority class.
We therefore treat \textbf{Bug-Fix F1} as the primary outcome metric,
as it reflects performance on the class of greatest practical interest,
and \textbf{MCC} as the primary aggregate metric, as it accounts for all
four confusion-matrix cells under class imbalance.
Non-Bug-Fix F1 is reported to assess whether improved Bug-Fix detection
comes at the cost of majority-class performance.

All methods are evaluated under 5-fold stratified cross-validation on
798 commits.
Unless stated otherwise, per-class metrics in
Tables~\ref{tab:rq1}--\ref{tab:confusion} pool predictions across all
folds before computing precision, recall, and F1.
Cross-fold stability (Table~\ref{tab:rq3-retrobug}) reports fold-wise
metrics with mean~$\pm$~standard deviation across the five splits.


\subsection{Evaluation Metrics}
\label{sec:evaluation-metrics}

We evaluate all methods on a binary classification task:
Bug-Fix (minority class) versus Non-Bug-Fix.
Predictions are compared against human-annotated labels using
standard per-class and aggregate metrics derived from the
confusion matrix (true positives~TP, true negatives~TN,
false positives~FP, false negatives~FN).

For each class $c \in \{\text{Non-Bug-Fix}, \text{Bug-Fix}\}$, we
report \emph{precision} ($\mathrm{TP}_c / (\mathrm{TP}_c + \mathrm{FP}_c)$),
\emph{recall} ($\mathrm{TP}_c / (\mathrm{TP}_c + \mathrm{FN}_c)$), and
\emph{F1} (harmonic mean of precision and recall).
We additionally report \emph{accuracy}
($(\mathrm{TP} + \mathrm{TN}) / N$) and the
\emph{Matthews Correlation Coefficient} (MCC)~\cite{matthews1975comparison}:
\[
  \mathrm{MCC} =
  \frac{\mathrm{TP}\cdot\mathrm{TN} - \mathrm{FP}\cdot\mathrm{FN}}
       {\sqrt{(\mathrm{TP}+\mathrm{FP})(\mathrm{TP}+\mathrm{FN})
              (\mathrm{TN}+\mathrm{FP})(\mathrm{TN}+\mathrm{FN})}}.
\]

Because the dataset is imbalanced (599 Non-Bug-Fix vs.\ 199 Bug-Fix
commits), accuracy alone can be misleading: a classifier that predicts
Non-Bug-Fix for nearly every commit achieves high accuracy while failing
on the minority class.
We therefore treat \textbf{Bug-Fix F1} as the primary outcome metric,
as it reflects performance on the class of greatest practical interest,
and \textbf{MCC} as the primary aggregate metric, as it accounts for all
four confusion-matrix cells under class imbalance.
Non-Bug-Fix F1 is reported to assess whether improved Bug-Fix detection
comes at the cost of majority-class performance.

All methods are evaluated under 5-fold stratified cross-validation on
798 commits.
Unless stated otherwise, per-class metrics in
Tables~\ref{tab:rq1}--\ref{tab:confusion} pool predictions across all
folds before computing precision, recall, and F1.
Cross-fold stability (Table~\ref{tab:rq3-retrobug}) reports fold-wise
metrics with mean~$\pm$~standard deviation across the five splits.


\subsection{Evaluation Metrics}
\label{sec:evaluation-metrics}

We evaluate all methods on a binary classification task:
Bug-Fix (minority class) versus Non-Bug-Fix.
Predictions are compared against human-annotated labels using
standard per-class and aggregate metrics derived from the
confusion matrix (true positives~TP, true negatives~TN,
false positives~FP, false negatives~FN).

For each class $c \in \{\text{Non-Bug-Fix}, \text{Bug-Fix}\}$, we
report \emph{precision} ($\mathrm{TP}_c / (\mathrm{TP}_c + \mathrm{FP}_c)$),
\emph{recall} ($\mathrm{TP}_c / (\mathrm{TP}_c + \mathrm{FN}_c)$), and
\emph{F1} (harmonic mean of precision and recall).
We additionally report \emph{accuracy}
($(\mathrm{TP} + \mathrm{TN}) / N$) and the
\emph{Matthews Correlation Coefficient} (MCC)~\cite{matthews1975comparison}:
\[
  \mathrm{MCC} =
  \frac{\mathrm{TP}\cdot\mathrm{TN} - \mathrm{FP}\cdot\mathrm{FN}}
       {\sqrt{(\mathrm{TP}+\mathrm{FP})(\mathrm{TP}+\mathrm{FN})
              (\mathrm{TN}+\mathrm{FP})(\mathrm{TN}+\mathrm{FN})}}.
\]

Because the dataset is imbalanced (599 Non-Bug-Fix vs.\ 199 Bug-Fix
commits), accuracy alone can be misleading: a classifier that predicts
Non-Bug-Fix for nearly every commit achieves high accuracy while failing
on the minority class.
We therefore treat \textbf{Bug-Fix F1} as the primary outcome metric,
as it reflects performance on the class of greatest practical interest,
and \textbf{MCC} as the primary aggregate metric, as it accounts for all
four confusion-matrix cells under class imbalance.
Non-Bug-Fix F1 is reported to assess whether improved Bug-Fix detection
comes at the cost of majority-class performance.

All methods are evaluated under 5-fold stratified cross-validation on
798 commits.
Unless stated otherwise, per-class metrics in
Tables~\ref{tab:rq1}--\ref{tab:confusion} pool predictions across all
folds before computing precision, recall, and F1.
Cross-fold stability (Table~\ref{tab:rq3-retrobug}) reports fold-wise
metrics with mean~$\pm$~standard deviation across the five splits.


\subsection{Evaluation Metrics}
\label{sec:evaluation-metrics}

We evaluate all methods on a binary classification task:
Bug-Fix (minority class) versus Non-Bug-Fix.
Predictions are compared against human-annotated labels using
standard per-class and aggregate metrics derived from the
confusion matrix (true positives~TP, true negatives~TN,
false positives~FP, false negatives~FN).

For each class $c \in \{\text{Non-Bug-Fix}, \text{Bug-Fix}\}$, we
report \emph{precision} ($\mathrm{TP}_c / (\mathrm{TP}_c + \mathrm{FP}_c)$),
\emph{recall} ($\mathrm{TP}_c / (\mathrm{TP}_c + \mathrm{FN}_c)$), and
\emph{F1} (harmonic mean of precision and recall).
We additionally report \emph{accuracy}
($(\mathrm{TP} + \mathrm{TN}) / N$) and the
\emph{Matthews Correlation Coefficient} (MCC)~\cite{matthews1975comparison}:
\[
  \mathrm{MCC} =
  \frac{\mathrm{TP}\cdot\mathrm{TN} - \mathrm{FP}\cdot\mathrm{FN}}
       {\sqrt{(\mathrm{TP}+\mathrm{FP})(\mathrm{TP}+\mathrm{FN})
              (\mathrm{TN}+\mathrm{FP})(\mathrm{TN}+\mathrm{FN})}}.
\]

Because the dataset is imbalanced (599 Non-Bug-Fix vs.\ 199 Bug-Fix
commits), accuracy alone can be misleading: a classifier that predicts
Non-Bug-Fix for nearly every commit achieves high accuracy while failing
on the minority class.
We therefore treat \textbf{Bug-Fix F1} as the primary outcome metric,
as it reflects performance on the class of greatest practical interest,
and \textbf{MCC} as the primary aggregate metric, as it accounts for all
four confusion-matrix cells under class imbalance.
Non-Bug-Fix F1 is reported to assess whether improved Bug-Fix detection
comes at the cost of majority-class performance.

All methods are evaluated under 5-fold stratified cross-validation on
798 commits.
Unless stated otherwise, per-class metrics in
Tables~\ref{tab:rq1}--\ref{tab:confusion} pool predictions across all
folds before computing precision, recall, and F1.
Cross-fold stability (Table~\ref{tab:rq3-retrobug}) reports fold-wise
metrics with mean~$\pm$~standard deviation across the five splits.
